\documentclass[12pt]{scrartcl}
\usepackage[hagavi]{azzam}

% Persamaan Diophantine dan Digit Bilangan - Materi
% Adaptasi dari 8 Outsource K Math/Pra-OSK/PersamaanDiophantine.tex,
% 8 Outsource K Math/Pra-OSP/PersamaanDiophantine.tex,
% 8 Outsource K Math/Pra-OSP/JumlahDigit.tex, dan
% 8 Outsource K Math/Pra-OSK/KoefisienBinomial.tex,
% disusun ulang mengikuti pola unit Materi/Latihan/Solusi.
% Gaya penulisan mengikuti mo_material.md: setiap teori langsung
% diikuti contoh soal, solusi contoh menumpuk persamaan ke bawah.
% Dari modul Koefisien Binomial hanya teorema binomial yang dipakai
% (sebagai alat modulo/keterbagian); identitas jumlahan Vandermonde
% dan stik hoki termasuk kombinatorika dan tidak masuk unit ini.
% Perbaikan dari modul sumber:
% - contoh x^2 - y! = 2001: solusi (25, 4) keliru, yang benar (45, 4);
% - contoh 3(xy+yz+zx) = 4xyz: tripel (1, 1, 6) keliru tulis,
%   yang benar (1, 6, 6).
% Dua contoh berat dari modul sumber (x^3 - y^3 = xy + 61 dan
% x^2 + y^2 = xy^2 + 25) dipindah ke Latihan agar contoh Materi
% tetap ringan.

\title{Persamaan Diophantine dan Digit Bilangan - Materi}
\date{21 Oktober 2026}

\begin{document}
\maketitle

\renewcommand*\contentsname{Daftar Isi}
\tableofcontents

\newpage

Modul ini merangkum dua tema teori bilangan yang sering muncul berdampingan
di OSK dan OSP: persamaan Diophantine dan sifat digit bilangan.
Enam seksi pertama membahas teknik standar persamaan Diophantine,
dari linear sampai keturunan tak hingga.
Dua seksi terakhir membahas jumlah digit dan teorema binomial sebagai alat bantu.
Setiap teori langsung diikuti contoh soal,
jadi bacalah teori lalu contohnya berpasangan.

\section{Persamaan Diophantine Linear}

Sebuah \istilah{persamaan Diophantine}{Diophantine equation} adalah persamaan
yang solusinya dicari dalam bilangan bulat.
Kadang daerah solusinya dipersempit lagi menjadi bilangan asli atau cacah,
sesuai redaksi soal.
Kasus paling dasar adalah persamaan linear dua variabel.

\begin{theorem}[Diophantine linear]
Misalkan $a$, $b$, dan $c$ bilangan bulat dengan $a$ dan $b$ tidak keduanya nol.
Tulis $d = FPB(a, b)$. Persamaan
\[ ax + by = c \]
mempunyai solusi bilangan bulat jika dan hanya jika $d \mid c$.
Apabila $(x_0, y_0)$ adalah satu solusi, semua solusinya adalah
\[ x = x_0 + \frac{b}{d}t, \qquad y = y_0 - \frac{a}{d}t \]
untuk $t$ bilangan bulat sebarang.
\end{theorem}

\begin{example}
Tentukan semua pasangan bilangan bulat $(x, y)$ yang memenuhi $6x + 10y = 4$.
\end{example}
\begin{proof}[Solusi]
Karena $FPB(6, 10) = 2$ membagi $4$, solusinya ada.
Bagi kedua ruas dengan $2$:
\[ 3x + 5y = 2. \]
Satu solusi mudah ditebak:
\[ x_0 = -1, \qquad y_0 = 1. \]
Solusi umumnya
\begin{align*}
x &= -1 + 5t, \\
y &= 1 - 3t,
\end{align*}
untuk setiap bilangan bulat $t$.
\end{proof}

\section{Faktorisasi}

Idenya sederhana. Ubah persamaan menjadi bentuk
\[ f_1 \cdot f_2 \cdots f_n = a \]
dengan $a$ konstanta, lalu periksa semua kemungkinan faktor dari $a$.
Tiga identitas berikut paling sering dipakai.

\begin{fact}[Identitas faktorisasi]
\begin{enumerate}
    \item \textbf{Simon's Favorite Factoring Trick.} Bentuk $xy + px + qy$
    selalu bisa dilengkapi menjadi $(x + q)(y + p)$ ditambah konstanta.
    \item \textbf{Sophie Germain.}
    \[ a^4 + 4b^4 = (a^2 - 2ab + 2b^2)(a^2 + 2ab + 2b^2). \]
    \item \textbf{Jumlah kubik.}
    \[ a^3 + b^3 + c^3 - 3abc = (a + b + c)(a^2 + b^2 + c^2 - ab - bc - ca). \]
\end{enumerate}
\end{fact}

\begin{example}
Tentukan semua pasangan bilangan bulat $(x, y)$ yang memenuhi $xy + 3x - 2y = 4$.
\end{example}
\begin{proof}[Solusi]
Lengkapi menjadi perkalian dua faktor:
\begin{align*}
(x - 2)(y + 3) &= xy + 3x - 2y - 6 \\
&= 4 - 6 \\
&= -2.
\end{align*}
Faktor dari $-2$ hanya $\pm 1$ dan $\pm 2$, jadi
\begin{align*}
x - 2 = 1, y + 3 = -2 &\quad \text{memberi} \quad (x, y) = (3, -5), \\
x - 2 = -1, y + 3 = 2 &\quad \text{memberi} \quad (x, y) = (1, -1), \\
x - 2 = 2, y + 3 = -1 &\quad \text{memberi} \quad (x, y) = (4, -4), \\
x - 2 = -2, y + 3 = 1 &\quad \text{memberi} \quad (x, y) = (0, -2).
\end{align*}
\end{proof}

\begin{example}
Tentukan semua bilangan asli $n$ sehingga $n^4 + 4$ merupakan bilangan prima.
\end{example}
\begin{proof}[Solusi]
Untuk $n = 1$ diperoleh $5$, yang prima.
Untuk $n$ genap, $n^4 + 4$ genap dan lebih dari $2$, jadi komposit.
Tinggal $n \ge 3$ ganjil. Faktorkan dengan Sophie Germain:
\[ n^4 + 4 = (n^2 - 2n + 2)(n^2 + 2n + 2). \]
Kedua faktornya lebih dari $1$, sebab
\[ n^2 - 2n + 2 = (n - 1)^2 + 1 \ge 2. \]
Jadi $n^4 + 4$ komposit. Satu-satunya jawaban $n = 1$.
\end{proof}

\begin{example}
Tentukan semua tripel bilangan asli $(x, y, z)$ yang memenuhi $x^3 + y^3 + z^3 = 3xyz$.
\end{example}
\begin{proof}[Solusi]
Pindahkan ruas dan pakai identitas jumlah kubik:
\[ (x + y + z)(x^2 + y^2 + z^2 - xy - yz - zx) = 0. \]
Faktor pertama positif. Faktor kedua bisa ditulis sebagai
\[ \frac{1}{2}\left((x - y)^2 + (y - z)^2 + (z - x)^2\right) = 0. \]
Jumlah tiga kuadrat nol hanya jika ketiganya nol:
\[ x = y = z. \]
Solusinya semua tripel $(t, t, t)$ dengan $t$ bilangan asli.
\end{proof}

\section{Ketaksamaan dan Bounding}

Kadang solusinya tidak bisa dicari langsung, tetapi ukurannya bisa dibatasi.
Setelah batasnya cukup kecil, sisa kasusnya tinggal diperiksa satu per satu.

\begin{fact}[Tiga cara bounding]
\begin{enumerate}
    \item \textbf{Mengurutkan.} Pada persamaan simetris, misalkan
    $x \le y \le z$ lalu ikat $x$ dari atas.
    \item \textbf{Ukuran pembagi.} Jika $a \mid b$ dengan $b \ne 0$, maka
    $|a| \le |b|$. Ketaksamaan ini sering memaksa variabelnya kecil.
    \item \textbf{Mengapit pangkat.} Di antara $m^n$ dan $(m+1)^n$ tidak ada
    pangkat $n$ lain. Jika suatu bentuk terjepit di antara dua kuadrat
    berurutan, ia bukan \istilah{bilangan kuadrat}{perfect square}.
\end{enumerate}
\end{fact}

\begin{example}
Tentukan semua pasangan bilangan asli $(x, y)$ dengan $x \mid y + 1$ dan $y \mid x + 1$.
\end{example}
\begin{proof}[Solusi]
Dari $x \mid y + 1$ diperoleh $x \le y + 1$,
dan dari $y \mid x + 1$ diperoleh $y \le x + 1$. Jadi
\[ |x - y| \le 1. \]
Jika $x = y$, maka $x \mid x + 1$ sehingga $x \mid 1$:
\[ (x, y) = (1, 1). \]
Jika $y = x + 1$, maka $x \mid x + 2$ sehingga $x \mid 2$:
\[ (x, y) = (1, 2) \quad \text{atau} \quad (2, 3). \]
Jika $x = y + 1$, simetris dengan kasus sebelumnya:
\[ (x, y) = (2, 1) \quad \text{atau} \quad (3, 2). \]
\end{proof}

\begin{example}[Pra-OSP Diophantine]
Carilah semua tripel bilangan asli $(x, y, z)$ yang memenuhi $3(xy + yz + zx) = 4xyz$.
\end{example}
\begin{proof}[Solusi]
Persamaannya simetris, jadi misalkan $x \le y \le z$.
Bagi kedua ruas dengan $3xyz$:
\[ \frac{1}{x} + \frac{1}{y} + \frac{1}{z} = \frac{4}{3}. \]
Karena $x \le y \le z$,
\[ \frac{4}{3} \le \frac{3}{x}, \qquad \text{sehingga} \qquad x \le 2. \]
Kasus $x = 1$ memberi
\[ \frac{1}{y} + \frac{1}{z} = \frac{1}{3} \le \frac{2}{y}, \]
sehingga $3 < y \le 6$. Pengecekan $y = 4, 5, 6$ memberi
\[ (x, y, z) = (1, 4, 12) \quad \text{dan} \quad (1, 6, 6), \]
beserta permutasinya. Kasus $x = 2$ memberi
\[ \frac{1}{y} + \frac{1}{z} = \frac{5}{6} \le \frac{2}{y}, \]
sehingga $2 \le y \le 2$. Nilai $y = 2$ memberi $z = 3$:
\[ (x, y, z) = (2, 2, 3), \]
beserta permutasinya.
\end{proof}

\begin{example}
Buktikan bahwa $n^2 + 3n + 1$ bukan bilangan kuadrat untuk setiap bilangan asli $n$.
\end{example}
\begin{proof}[Solusi]
Apit bentuk itu di antara dua kuadrat berurutan:
\begin{align*}
(n + 1)^2 &= n^2 + 2n + 1 < n^2 + 3n + 1, \\
(n + 2)^2 &= n^2 + 4n + 4 > n^2 + 3n + 1.
\end{align*}
Karena letaknya strikt di antara dua kuadrat berurutan,
$n^2 + 3n + 1$ bukan bilangan kuadrat.
\end{proof}

\section{Modulo}

Analisis modulo sering langsung menghabisi soal.
Pilih modulus yang membuat residu pangkatnya sedikit,
lalu tunjukkan persamaannya mustahil atau sisakan sedikit kasus untuk diperiksa.
Daftar di bawah ini sengaja ringkas.
Versi lengkapnya, beserta teorema Euler, Fermat Kecil, dan Wilson,
ada di unit Modulo dan Keterbagian.

\begin{fact}[Residu pangkat]
\begin{enumerate}
    \item Sisa $x^2$ dibagi $4$ hanya bisa $0$ atau $1$.
    Sisa $x^2$ dibagi $8$ hanya bisa $0$, $1$, atau $4$.
    \item Sisa $x^3$ dibagi $9$ hanya bisa $0$, $1$, atau $8$.
    \item Sisa $x^4$ dibagi $5$ hanya bisa $0$ atau $1$.
\end{enumerate}
\end{fact}

\begin{example}
Buktikan tidak ada bilangan bulat $x$ dan $y$ dengan $x^2 + y^2 = 2023$.
\end{example}
\begin{proof}[Solusi]
Tinjau persamaan itu modulo $4$:
\[ x^2 + y^2 \equiv 2023 \equiv 3 \pmod{4}. \]
Sisa kuadrat modulo $4$ hanya $0$ atau $1$,
sehingga jumlah dua sisa hanya bisa $0$, $1$, atau $2$.
Nilai $3$ tidak pernah muncul.
\end{proof}

\begin{example}[Pra-OSP Diophantine]
Carilah semua pasangan bilangan asli $(x, y)$ yang memenuhi $x^2 - y! = 2001$.
\end{example}
\begin{proof}[Solusi]
Untuk $y \ge 3$, nilai $y! + 2001$ habis dibagi $3$,
sehingga $3 \mid x$ dan $x^2 \equiv 0 \pmod{9}$.
Karena $2001 \equiv 3 \pmod{9}$,
\[ y! = x^2 - 2001 \equiv -3 \equiv 6 \pmod{9}. \]
Tetapi $y!$ habis dibagi $9$ untuk $y \ge 6$.
Jadi $y \le 5$, dan pengecekan langsung memberi
\begin{align*}
y = 1 &: x^2 = 2002, \quad 44^2 < 2002 < 45^2, \\
y = 2 &: x^2 = 2003 \equiv 2 \pmod{3}, \\
y = 3 &: x^2 = 2007, \quad 44^2 < 2007 < 45^2, \\
y = 4 &: x^2 = 2025 = 45^2, \\
y = 5 &: x^2 = 2121, \quad 46^2 < 2121 < 47^2.
\end{align*}
Satu-satunya solusi $(x, y) = (45, 4)$.
\end{proof}

\section{Keturunan Tak Hingga}

Metode \istilah{keturunan tak hingga}{infinite descent} adalah senjata standar
untuk persamaan homogen berderajat tinggi. Idenya satu baris.

\begin{fact}[Prinsip keturunan tak hingga]
Tidak ada barisan bilangan asli tak hingga yang turun tegas.
Akibatnya, jika setiap solusi tak trivial dari suatu persamaan melahirkan
solusi tak trivial yang lebih kecil, persamaan itu tidak punya solusi tak trivial.
\end{fact}

\begin{example}
Tentukan semua solusi bilangan bulat dari $x^2 + y^2 = 3z^2$.
\end{example}
\begin{proof}[Solusi]
Jelas $(0, 0, 0)$ memenuhi. Andaikan ada solusi tak nol.
Modulo $3$, sisa kuadrat hanya $0$ atau $1$, sehingga
\[ x^2 + y^2 \equiv 0 \pmod{3} \]
memaksa $x$ dan $y$ habis dibagi $3$.
Tulis $x = 3x_1$ dan $y = 3y_1$:
\[ 9x_1^2 + 9y_1^2 = 3z^2, \qquad \text{yaitu} \qquad z^2 = 3(x_1^2 + y_1^2). \]
Jadi $3 \mid z$ juga. Tulis $z = 3z_1$:
\[ x_1^2 + y_1^2 = 3z_1^2. \]
Tripel $(x_1, y_1, z_1)$ solusi yang lebih kecil:
\[ |x_1| + |y_1| + |z_1| = \frac{|x| + |y| + |z|}{3}. \]
Pengulangan tanpa henti mustahil untuk bilangan cacah,
jadi satu-satunya solusi $(x, y, z) = (0, 0, 0)$.
\end{proof}

\begin{example}[Pra-OSP Diophantine]
Carilah semua solusi bilangan bulat dari $x^4 + y^4 + z^4 = 9u^4$.
\end{example}
\begin{proof}[Solusi]
Pangkat $4$ menyarankan modulo $5$, karena $t^4 \equiv 0$ atau $1 \pmod{5}$.
Jelas $(0, 0, 0, 0)$ memenuhi. Jika $u = 0$, maka $x = y = z = 0$ juga.
Andaikan ada solusi dengan $u \ne 0$.
Jika $5 \nmid u$, ruas kanan kongruen $9 \equiv 4 \pmod{5}$,
sementara ruas kiri, jumlah tiga nilai $0$ atau $1$, paling besar $3$.
Kontradiksi. Jadi $5 \mid u$, dan dari
\[ x^4 + y^4 + z^4 = 9u^4 \equiv 0 \pmod{5} \]
diperoleh $5 \mid x$, $5 \mid y$, dan $5 \mid z$.
Tulis $x = 5x_1$ dan seterusnya, lalu bagi kedua ruas dengan $5^4$:
\[ x_1^4 + y_1^4 + z_1^4 = 9u_1^4. \]
Solusi baru ini lebih kecil, sebab $|u_1| = |u|/5 < |u|$.
Pengulangan tanpa henti mustahil, jadi satu-satunya solusi
\[ (x, y, z, u) = (0, 0, 0, 0). \]
\end{proof}

\section{Diskriminan}

Pandang persamaan sebagai kuadrat dalam satu variabel.
Akarnya bulat hanya jika \istilah{diskriminan}{discriminant}nya bilangan kuadrat.

\begin{fact}[Syarat diskriminan]
Persamaan kuadrat $at^2 + bt + c = 0$ dengan $a$, $b$, $c$ bilangan bulat
mempunyai akar rasional hanya jika $b^2 - 4ac$ adalah bilangan kuadrat.
\end{fact}

\begin{example}
Tentukan semua pasangan bilangan bulat $(x, y)$ yang memenuhi $y^2 = x^2 + x + 1$.
\end{example}
\begin{proof}[Solusi]
Pandang sebagai kuadrat dalam $x$:
\[ x^2 + x + (1 - y^2) = 0. \]
Diskriminannya harus bilangan kuadrat:
\[ D = 1 - 4(1 - y^2) = 4y^2 - 3 = t^2 \]
untuk suatu bilangan cacah $t$. Faktorkan:
\[ (2y - t)(2y + t) = 3. \]
Kedua faktornya ganjil, dan $3 = 1 \cdot 3 = (-1) \cdot (-3)$.
Pasangan $(1, 3)$ dan $(3, 1)$ memberi $y = 1$,
pasangan $(-1, -3)$ dan $(-3, -1)$ memberi $y = -1$.
Dari $x = (-1 \pm t)/2$ dengan $t = 1$ diperoleh
\[ (x, y) \in \{(0, 1), (0, -1), (-1, 1), (-1, -1)\}. \]
\end{proof}

\section{Jumlah Digit}

\begin{definition}
Untuk bilangan asli $n$, tulis $S(n)$ untuk \istilah{jumlah digit}{digit sum}
dari $n$, yaitu hasil penjumlahan semua digit desimalnya.
\end{definition}

\begin{fact}[Sifat jumlah digit]
Untuk bilangan asli $a$, $m$, dan $n$ berlaku:
\begin{enumerate}
    \item $S(an) \le aS(n)$.
    \item $S(n) \equiv n \pmod{9}$.
    \item $S(mn) \le S(m)S(n)$.
\end{enumerate}
\end{fact}

\begin{example}
Tentukan semua bilangan asli dua digit $n$ dengan $n = 2S(n)$.
\end{example}
\begin{proof}[Solusi]
Tulis $n = 10a + b$ dengan $a$ dan $b$ digit, $a \ne 0$:
\begin{align*}
10a + b &= 2(a + b), \\
8a &= b.
\end{align*}
Karena $b$ satu digit, $a = 1$ dan $b = 8$.
Jadi $n = 18$.
\end{proof}

\begin{example}
Tentukan semua bilangan asli $n$ dengan $n - S(n) = 45$.
\end{example}
\begin{proof}[Solusi]
Dari $n - S(n) = 45$ diperoleh $45 \le n < 100$,
sebab $n - S(n) \ge 73 > 45$ untuk $n \ge 100$.
Jadi $n$ dua digit. Tulis $n = 10a + b$:
\begin{align*}
10a + b - (a + b) &= 45, \\
9a &= 45, \\
a &= 5.
\end{align*}
Digit $b$ bebas. Solusinya $n \in \{50, 51, \dots, 59\}$.
\end{proof}

\begin{example}
Tentukan semua bilangan asli $n$ dengan $n + S(n) = 100$.
\end{example}
\begin{proof}[Solusi]
Karena $S(n) \ge 1$, berlaku $n \le 99$, sehingga $S(n) \le 18$ dan $n \ge 82$.
Jadi $n$ dua digit. Tulis $n = 10a + b$:
\begin{align*}
10a + b + a + b &= 100, \\
11a + 2b &= 100.
\end{align*}
Nilai $2b = 100 - 11a$ harus genap di antara $0$ dan $18$,
yang hanya terjadi untuk $a = 8$ dan $b = 6$.
Jadi $n = 86$.
\end{proof}

\begin{remark}
Bukti ketiga sifat jumlah digit di atas diminta sebagai soal latihan pertama.
\end{remark}

\section{Teorema Binomial}

Koefisien binomial muncul di banyak tempat.
Di unit ini ia hanya alat bantu:
menguraikan $(a + b)^n$ untuk argumen modulo dan keterbagian.

\begin{definition}
Untuk bilangan cacah $n$ dan $r$ dengan $r \le n$,
\istilah{koefisien binomial}{binomial coefficient} $\binom{n}{r}$ adalah
\[ \binom{n}{r} = \frac{n!}{r!(n-r)!}, \]
banyak cara memilih $r$ benda dari $n$ benda.
\end{definition}

Hasil utamanya adalah \istilah{teorema binomial}{binomial theorem}.

\begin{theorem}[Teorema binomial]
Untuk setiap bilangan cacah $n$ dan bilangan $x$, $y$ sebarang berlaku
\[ (x + y)^n = \sum_{r=0}^{n} \binom{n}{r} x^{n-r} y^r. \]
\end{theorem}

\begin{example}
Buktikan bahwa $49$ membagi $8^7 - 1$.
\end{example}
\begin{proof}[Solusi]
Tulis $8 = 1 + 7$ lalu uraikan:
\begin{align*}
8^7 - 1 &= (1 + 7)^7 - 1 \\
&= \binom{7}{1} \cdot 7 + \binom{7}{2} \cdot 7^2 + \dots + 7^7.
\end{align*}
Suku pertama $\binom{7}{1} \cdot 7 = 49$.
Semua suku berikutnya memuat $7^2$ atau lebih tinggi.
Jadi $49$ membagi jumlahnya.
\end{proof}

\begin{example}
Tentukan sisa pembagian $4^{10}$ oleh $9$.
\end{example}
\begin{proof}[Solusi]
Tulis $4 = 1 + 3$ lalu uraikan:
\begin{align*}
4^{10} = (1 + 3)^{10} &= 1 + \binom{10}{1} \cdot 3 + \binom{10}{2} \cdot 3^2 + \dots \\
&= 1 + 30 + 9 \cdot (\text{bilangan bulat}).
\end{align*}
Modulo $9$ suku terakhir lenyap:
\[ 4^{10} \equiv 1 + 30 \equiv 4 \pmod{9}. \]
Sisanya $4$.
\end{proof}

\section{Referensi}
\begin{itemize}
    \item Modul Pra-OSK: Persamaan Diophantine (Azzam dan Akeyla).
    \item Modul Pra-OSP: Persamaan Diophantine (Azzam dan Akeyla).
    \item Modul Pra-OSP: Jumlah Digit (Azzam dan Akeyla).
    \item Modul Pra-OSK: Koefisien Binomial (Azzam dan Akeyla).
\end{itemize}

\end{document}
